\chapter{Vormen en ruimte}\label{ch:g1:shapes}

Ramen zijn rechthoeken, wielen zijn cirkels, in daken zitten
driehoeken verstopt: vormen zijn overal. In dit hoofdstuk leer je ze
herkennen door zijden en hoeken te tellen, en leer je precies zeggen
\emph{waar} iets is: links, rechts, boven, onder, tussen.

\section{De vier basisvormen}

\begin{definition}[Vierkant, rechthoek, driehoek, cirkel]\label{def:g1:shapes:def}
\begin{itemize}
  \item de \emph{driehoek}\index{driehoek} heeft $3$ zijden en $3$ hoeken;
  \item het \emph{vierkant}\index{vierkant} heeft $4$ gelijke zijden en $4$ hoeken;
  \item de \emph{rechthoek}\index{rechthoek} heeft $4$ zijden en $4$ hoeken, waarbij de
        zijden die tegenover elkaar liggen gelijk zijn (een uitgerekt
        vierkant);
  \item de \emph{cirkel}\index{cirkel} is helemaal rond: geen zijden, geen hoeken.
\end{itemize}
\end{definition}

\begin{omfigure}
\begin{tikzpicture}[scale=0.7]
  \draw[thick, fill=omDef!12] (0,0) -- (2,0) -- (1,1.7) -- cycle;
  \node[below, align=center, font=\small] at (1,-0.25)
    {driehoek:\\$3$ zijden};
  \begin{scope}[xshift=3.6cm]
  \draw[thick, fill=omThm!12] (0,0) rectangle (1.7,1.7);
  \node[below, align=center, font=\small] at (0.85,-0.25)
    {vierkant:\\$4$ gelijke zijden};
  \end{scope}
  \begin{scope}[xshift=6.9cm]
  \draw[thick, fill=omMeth!12] (0,0.05) rectangle (2.8,1.65);
  \node[below, align=center, font=\small] at (1.4,-0.25)
    {rechthoek:\\$4$ zijden};
  \end{scope}
  \begin{scope}[xshift=11.3cm]
  \draw[thick, fill=omProp!12] (0.9,0.85) circle (0.9);
  \node[below, align=center, font=\small] at (0.9,-0.25)
    {cirkel:\\geen hoek};
  \end{scope}
\end{tikzpicture}

{\small De vier basisvormen. Vertrouw bij het benoemen niet op de
eerste blik: \emph{tel} de zijden en de hoeken.}
\end{omfigure}

\begin{example}[Gekantelde vormen blijven vormen]\label{ex:g1:shapes:tilted}
Een vierkant dat op zijn punt staat, blijft een vierkant --- vier
gelijke zijden, vier hoeken --- ook al lijkt het op een vlieger. Een
vorm draaien verandert niet wat het is.
\end{example}

\begin{example}[Vormenjacht]\label{ex:g1:shapes:hunt}
In de klas: het bord is een rechthoek, de klok een cirkel, een
gevouwen servet kan een driehoek zijn, sommige ramen zijn vierkanten.
Echte voorwerpen zijn geen perfecte vormen, maar hun vlakken komen er
dichtbij (lichamen en hun zijvlakken komen terug in
\cref{ch:g2:shapes}).
\end{example}

\section{Waar is het?}

\begin{definition}[Plaatswoorden]\label{def:g1:shapes:position}
Om te zeggen waar iets is: \emph{links} / \emph{rechts},
\emph{boven} / \emph{onder}, \emph{op} / \emph{onder},
\emph{voor} / \emph{achter}, \emph{tussen},
\emph{binnen} / \emph{buiten}.
\end{definition}

\begin{omfigure}
\begin{tikzpicture}[scale=0.75]
  \draw[thick, fill=omDef!12] (0,0) circle (0.55);
  \draw[thick, fill=omThm!12] (1.9,-0.55) rectangle (3.0,0.55);
  \draw[thick, fill=omMeth!12] (4.4,-0.5) -- (5.5,-0.5) -- (4.95,0.45)
    -- cycle;
  \node[below, font=\small] at (0,-0.75) {cirkel};
  \node[below, font=\small] at (2.45,-0.75) {vierkant};
  \node[below, font=\small] at (4.95,-0.75) {driehoek};
  \node[right, font=\small, align=left] at (6.3,0)
    {het vierkant ligt \emph{tussen}\\de cirkel en de driehoek;\\
    de cirkel ligt \emph{links} van het vierkant};
\end{tikzpicture}

{\small Plaatswoorden in actie. Let op: \emph{jouw} links en de links
van iemand die je aankijkt zijn tegengesteld!}
\end{omfigure}

\section{Paden op een raster}

\begin{method}[Een pad beschrijven]\label{met:g1:shapes:path}
Op ruitjespapier is een pad een rij stappen:
$\to$ (één vakje naar rechts), $\leftarrow$ (naar links), $\uparrow$
(naar boven), $\downarrow$ (naar beneden).
\begin{enumerate}
  \item Een pad volgen: begin op het gemarkeerde vakje en volg de
        pijlen één voor één;
  \item een pad beschrijven: loop het in gedachten en schrijf de
        pijlen op volgorde;
  \item twee verschillende paden kunnen dezelfde twee vakjes
        verbinden --- het ene korter, het andere langer.
\end{enumerate}
\end{method}

\begin{omfigure}
\begin{tikzpicture}[scale=0.6]
  \draw[gray!40] (0,0) grid (7,4);
  \fill[omDef!40] (0.5,0.5) +(-0.42,-0.42) rectangle +(0.42,0.42);
  \node[font=\footnotesize] at (0.5,0.5) {start};
  \fill[omThm!40] (5.5,2.5) +(-0.42,-0.42) rectangle +(0.42,0.42);
  \node[font=\footnotesize] at (5.5,2.5) {doel};
  % zeven pijlen van één vakje, met kleine tussenruimtes
  \foreach \a/\b in {0.5/1.5, 1.5/2.5}
    \draw[-Stealth, omDef, very thick] (\a+0.1,0.5) -- (\b-0.1,0.5);
  \foreach \a/\b in {0.5/1.5, 1.5/2.5}
    \draw[-Stealth, omDef, very thick] (2.5,\a+0.1) -- (2.5,\b-0.1);
  \foreach \a/\b in {2.5/3.5, 3.5/4.5}
    \draw[-Stealth, omDef, very thick] (\a+0.1,2.5) -- (\b-0.1,2.5);
  \draw[-Stealth, omDef, very thick] (4.6,2.5) -- (5.05,2.5);
\end{tikzpicture}

{\small Een pad van start naar doel, één pijl per stap:
$\to\ \to\ \uparrow\ \uparrow\ \to\ \to\ \to$ --- zeven stappen. Kun
jij een ander pad met even veel stappen vinden?}
\end{omfigure}

\section{Oefeningen}

\begin{exercise}[$\star$]\label{exo:g1:shapes:1}
Tel bij elke vorm de zijden en de hoeken en noem de vorm: een vorm met
$3$ hoeken; een ronde vorm zonder hoek; een vorm met $4$ gelijke
zijden.
\end{exercise}

\begin{exercise}[$\star$]\label{exo:g1:shapes:2}
Zoek thuis: twee rechthoeken, één cirkel, één driehoek, één vierkant.
(Deuren? borden? verkeersborden?)
\end{exercise}

\begin{exercise}[$\star$]\label{exo:g1:shapes:3}
Waar of niet waar? ``Een vierkant dat op zijn punt staat, is geen
vierkant meer.'' Leg uit met \cref{ex:g1:shapes:tilted}.
\end{exercise}

\begin{exercise}[$\star$]\label{exo:g1:shapes:4}
Teken op ruitjespapier: een vierkant van $3$ bij $3$ vakjes; een
rechthoek van $5$ bij $2$; een driehoek met een diagonaal van het
raster als zijde.
\end{exercise}

\begin{exercise}[$\star$]\label{exo:g1:shapes:5}
Zet drie stukken speelgoed op een rij. Beschrijf de rij met
plaatswoorden: wat staat tussen de andere twee? Wat staat links?
\end{exercise}

\begin{exercise}[$\star$]\label{exo:g1:shapes:6}
Teken een cirkel. Teken een kruisje \emph{binnen} de cirkel, een
sterretje \emph{buiten} de cirkel, en een stip \emph{op} de cirkel
zelf.
\end{exercise}

\begin{exercise}[$\star$]\label{exo:g1:shapes:7}
Markeer op ruitjespapier een startvakje. Volg:
$\to\ \uparrow\ \to\ \uparrow\ \to$. Markeer waar je aankomt. Schrijf
daarna het pad dat rechtstreeks terug naar de start gaat.
\end{exercise}

\begin{exercise}[$\star$]\label{exo:g1:shapes:8}
Hoeveel hoeken zijn er in totaal bij twee driehoeken en één vierkant?
En bij één cirkel en twee rechthoeken?
\end{exercise}

\begin{exercise}[$\star\star$]\label{exo:g1:shapes:9}
Neem $12$ luciferstokjes. Kun je met \emph{alle} stokjes een vierkant
leggen, met hele stokjes en even veel op elke zijde? Hoeveel stokjes
per zijde? En zou het met $10$ stokjes lukken?
\end{exercise}
